% Macros shared by all chapters of the course "Variational Methods".
% Loaded by both handoutsbeamer.tex (article mode) and overheadsbeamer.tex (presentation mode).

% A figure that is shown full-frame in the slides and as a captioned float in the handout.
%   \figslide[<includegraphics options for the slide>]{<file>}{<caption>}{<label>}
% In the handout, the figure is always set at 0.92 of the line width.
\newcommand{\figslide}[4][height=0.62\textheight]{%
  \onp{\begin{center}\includegraphics[#1]{#2}\end{center}\vspace{-0.6em}{\tiny #3}}%
  \ona{\begin{figure}[htb]\centering\includegraphics[width=0.92\linewidth,height=0.42\textheight,keepaspectratio]{#2}\caption{#3}\label{#4}\end{figure}}}

% Cross references to chapters and lectures
% "Chapter" in the handout, "Lecture" in the slides
\newcommand{\Chap}{\only<article>{Chapter}\only<presentation>{Lecture}}
\newcommand{\Chaps}{\only<article>{Chapters}\only<presentation>{Lectures}}
\newcommand{\chap}{\only<article>{chapter}\only<presentation>{lecture}}
\newcommand{\chaps}{\only<article>{chapters}\only<presentation>{lectures}}

% Notation
% (\bs is the backslash in declarns.tex; use \boldsymbol directly)
\newcommand{\Herm}{\mathrm{Herm}}
\DeclareMathOperator*{\argmin}{arg\,min}
\DeclareMathOperator*{\argmax}{arg\,max}
\renewcommand*{\Tr}{\ensuremath{\operatorname{Tr}}}

\DeclareMathOperator{\opt}{opt}
\DeclareMathOperator{\dist}{dist}
\DeclareMathOperator{\spn}{span}

\providecommand{\R}{\mathbb{R}}
\providecommand{\C}{\mathbb{C}}
\providecommand{\N}{\mathbb{N}}
\usepackage{dsfont}
\providecommand{\one}{}\renewcommand{\one}{\mathds{1}}   % identity operator (dsfont has a proper double-struck 1; \mathbb{1} has none)
\newcommand{\todo}[1]{\textcolor{red}{[TODO: #1]}}

% Chapter cross-references that also work when a single chapter is compiled on its own:
% \chref{C.xxx} expands to \ref{C.xxx} if the label is known, and to the fixed chapter
% number of the course plan otherwise.
\makeatletter
\@namedef{chapnum@C.Motivation}{1}
\@namedef{chapnum@C.EQA}{2}
\@namedef{chapnum@C.Inapprox}{3}
\@namedef{chapnum@C.Annealing}{4}
\@namedef{chapnum@C.VQE}{5}
\@namedef{chapnum@C.QAOA}{6}
\@namedef{chapnum@C.PSR}{7}
\@namedef{chapnum@C.IterationComplexity}{8}
\@namedef{chapnum@C.PerIteration}{9}
\@namedef{chapnum@C.WSRounded}{10}
\@namedef{chapnum@C.WSContinuous}{11}
\@namedef{chapnum@C.VQLS}{12}
\@namedef{chapnum@C.Portfolio}{13}
\newcommand{\chref}[1]{\ifcsname r@#1\endcsname\ref{#1}\else\ifcsname chapnum@#1\endcsname\csname chapnum@#1\endcsname\else\ref{#1}\fi\fi}
\makeatother

% Code listings (Python) shown in both the slides and the handout
\lstset{language=Python,basicstyle=\ttfamily\scriptsize,keywordstyle=\color{blue!60!black}\bfseries,commentstyle=\color{green!40!black}\itshape,stringstyle=\color{red!50!black},showstringspaces=false,columns=fullflexible,keepspaces=true,breaklines=true,frame=single,framerule=0.3pt,rulecolor=\color{gray},xleftmargin=1em,xrightmargin=1em,numbers=none}

% TikZ figures drawn in the chapters
\usepackage{pgfplots}
\pgfplotsset{compat=1.18}
\usetikzlibrary{calc,arrows.meta,positioning,backgrounds}

% Angstrom that works in both text and math, in both modes
\providecommand{\angstrom}{\ensuremath{\mathring{\mathrm{A}}}}   % \r is redefined by the template's macros, so use the math ring accent
